\nonstopmode{}
\documentclass[a4paper]{book}
\usepackage[times,inconsolata,hyper]{Rd}
\usepackage{makeidx}
\makeatletter\@ifl@t@r\fmtversion{2018/04/01}{}{\usepackage[utf8]{inputenc}}\makeatother
% \usepackage{graphicx} % @USE GRAPHICX@
\makeindex{}
\begin{document}
\chapter*{}
\begin{center}
{\textbf{\huge Package `helpeRs'}}
\par\bigskip{\large \today}
\end{center}
\ifthenelse{\boolean{Rd@use@hyper}}{\hypersetup{pdftitle = {helpeRs: Helper Functions for High-Quality Presentation of Quantitative Research}}}{}
\ifthenelse{\boolean{Rd@use@hyper}}{\hypersetup{pdfauthor = {Connor Jerzak}}}{}
\begin{description}
\raggedright{}
\item[Title]\AsIs{Helper Functions for High-Quality Presentation of Quantitative Research}
\item[Version]\AsIs{0.0}
\item[Description]\AsIs{Helper functions for high-quality presentation of quantitative
research. Provides tools for generating publication-ready regression tables
in LaTeX format, with support for clustered standard errors, bootstrap
inference, and instrumental variables. Also includes visualization utilities
for creating heat maps and other graphics.}
\item[Depends]\AsIs{R (>= 3.3.3)}
\item[License]\AsIs{GPL-3}
\item[Encoding]\AsIs{UTF-8}
\item[Imports]\AsIs{fields,
ggExtra,
ggplot2,
interp,
lmtest,
missForest,
plyr,
sandwich,
stargazer,
stringr,
viridis}
\item[RoxygenNote]\AsIs{7.3.3}
\item[Suggests]\AsIs{MASS,
testthat (>= 3.0.0)}
\item[Config/testthat/edition]\AsIs{3}
\end{description}
\Rdcontents{Contents}
\HeaderA{cols2numeric}{Convert data frame columns to numeric when possible}{cols2numeric}
%
\begin{Description}
Attempts to coerce each column of a data frame to numeric using
\code{\LinkA{f2n}{f2n}}. Columns that cannot be coerced (i.e., result in all
\code{NA} values) are left unchanged. This is useful for cleaning data
where numeric values may be stored as factors or character strings.
\end{Description}
%
\begin{Usage}
\begin{verbatim}
cols2numeric(x)
\end{verbatim}
\end{Usage}
%
\begin{Arguments}
\begin{ldescription}
\item[\code{x}] A data frame to process.
\end{ldescription}
\end{Arguments}
%
\begin{Value}
A data frame with the same dimensions as \code{x}. Columns that
could be converted to numeric are returned as numeric; others retain
their original type.
\end{Value}
%
\begin{SeeAlso}
\code{\LinkA{f2n}{f2n}} for converting individual vectors,
\code{\LinkA{colSummmary}{colSummmary}} for summarizing columns
\end{SeeAlso}
%
\begin{Examples}
\begin{ExampleCode}
# Data frame with mixed types
df <- data.frame(
  a = factor(c("1.5", "2.5", "3.5")),
  b = c("x", "y", "z"),
  c = c(1, 2, 3),
  stringsAsFactors = FALSE
)

# Convert what can be converted
df_numeric <- cols2numeric(df)
sapply(df_numeric, class)
# a becomes numeric, b stays character, c stays numeric

\end{ExampleCode}
\end{Examples}
\HeaderA{colSummmary}{Summarize each column of a data frame}{colSummmary}
%
\begin{Description}
Computes a single summary value for each column of a data frame. Numeric
columns are summarized by their mean, while non-numeric columns are
summarized by their mode (most frequent value). This is primarily used when
constructing prediction grids for visualization, where you need representative
values for variables not being varied.
\end{Description}
%
\begin{Usage}
\begin{verbatim}
colSummmary(x)
\end{verbatim}
\end{Usage}
%
\begin{Arguments}
\begin{ldescription}
\item[\code{x}] A data frame to summarize.
\end{ldescription}
\end{Arguments}
%
\begin{Details}
The function uses \code{\LinkA{f2n}{f2n}} internally to attempt numeric conversion,
so factor columns with numeric-looking levels will be treated as numeric.
\end{Details}
%
\begin{Value}
A named vector with one element per column of \code{x}. Names
correspond to column names. Numeric summaries are means; non-numeric
summaries are the most frequent value.
\end{Value}
%
\begin{SeeAlso}
\code{\LinkA{cols2numeric}{cols2numeric}} for converting columns to numeric,
\code{\LinkA{MakeHeatMap}{MakeHeatMap}} which uses this function internally
\end{SeeAlso}
%
\begin{Examples}
\begin{ExampleCode}
# Summarize mtcars
colSummmary(mtcars[, 1:4])

# Mixed numeric and character columns
df <- data.frame(
  x = c(1, 2, 3, 4),
  y = c("a", "b", "a", "a")
)
colSummmary(df)  # Returns mean of x and mode of y

\end{ExampleCode}
\end{Examples}
\HeaderA{f2n}{Convert factors to numeric values}{f2n}
%
\begin{Description}
Converts a vector to numeric by first coercing to character. This is the
correct way to extract numeric values from factors, as direct coercion with
\code{as.numeric()} returns factor level indices rather than the underlying
values.
\end{Description}
%
\begin{Usage}
\begin{verbatim}
f2n(x)
\end{verbatim}
\end{Usage}
%
\begin{Arguments}
\begin{ldescription}
\item[\code{x}] A vector to convert, typically a factor whose levels are numeric
strings (e.g., \code{factor(c("1.5", "2.3", "3.1"))}).
\end{ldescription}
\end{Arguments}
%
\begin{Value}
A numeric vector. Values that cannot be coerced to numeric become
\code{NA} with a warning.
\end{Value}
%
\begin{SeeAlso}
\code{\LinkA{cols2numeric}{cols2numeric}} for converting multiple columns of a
data frame
\end{SeeAlso}
%
\begin{Examples}
\begin{ExampleCode}
# Factor with numeric levels
x <- factor(c("1.5", "2.3", "3.1"))

# Wrong way - returns level indices (1, 2, 3)
as.numeric(x)

# Correct way - returns actual values
f2n(x)  # Returns c(1.5, 2.3, 3.1)

\end{ExampleCode}
\end{Examples}
\HeaderA{fixZeroEndings}{Ensure numbers have a fixed number of decimal places}{fixZeroEndings}
%
\begin{Description}
Pads numeric strings with trailing zeros so that all values have exactly
\code{roundAt} digits after the decimal point. This ensures consistent
formatting in regression tables where column alignment matters.
\end{Description}
%
\begin{Usage}
\begin{verbatim}
fixZeroEndings(zr, roundAt = 2)
\end{verbatim}
\end{Usage}
%
\begin{Arguments}
\begin{ldescription}
\item[\code{zr}] A character or numeric vector to process. Numeric values are first
converted to character.

\item[\code{roundAt}] Integer specifying the desired number of decimal places.
Default is 2.
\end{ldescription}
\end{Arguments}
%
\begin{Value}
A character vector with values padded to have exactly \code{roundAt}
decimal places. Values without a decimal point receive one followed by
the appropriate number of zeros.
\end{Value}
%
\begin{SeeAlso}
\code{\LinkA{GetTableEntry}{GetTableEntry}} which uses this function to format
coefficient estimates
\end{SeeAlso}
%
\begin{Examples}
\begin{ExampleCode}
# Pad to 2 decimal places (default)
fixZeroEndings(c("1.5", "2", "3.14"))
# Returns: c("1.50", "2.00", "3.14")

# Pad to 3 decimal places
fixZeroEndings(c(1.5, 2.33, 3), roundAt = 3)
# Returns: c("1.500", "2.330", "3.000")

\end{ExampleCode}
\end{Examples}
\HeaderA{FullTransformer}{Clean and reorder regression tables for publication}{FullTransformer}
%
\begin{Description}
Takes a matrix of regression output and performs string replacements,
row name cleaning, and reordering so that the resulting table is suitable
for inclusion in LaTeX documents. This function is typically called internally
by \code{\LinkA{Tables2Tex}{Tables2Tex}} but can be used directly for custom table processing.
\end{Description}
%
\begin{Usage}
\begin{verbatim}
FullTransformer(t_FULL, COLNAMES_VEC, name_conversion_matrix = NULL)
\end{verbatim}
\end{Usage}
%
\begin{Arguments}
\begin{ldescription}
\item[\code{t\_FULL}] Matrix or data frame containing the raw regression table output,
typically from \code{\LinkA{GetTableEntry}{GetTableEntry}}.

\item[\code{COLNAMES\_VEC}] Character vector of column names to apply to the final table.
Length must match the number of columns in \code{t\_FULL}.

\item[\code{name\_conversion\_matrix}] Optional two-column matrix for renaming row labels.
Column 1 contains regex patterns to match against row names,
column 2 contains the replacement names. If \code{NULL} (default), no
renaming is performed.
\end{ldescription}
\end{Arguments}
%
\begin{Details}
The function performs several transformations:
\begin{itemize}

\item{} Renames rows using the optional \code{name\_conversion\_matrix}
\item{} Moves summary statistics (R-squared, observations, etc.) to the bottom
\item{} Expands camelCase variable names into readable format
\item{} Handles \code{as.factor()} variable names by extracting factor levels
\item{} Preserves common acronyms (GDP, AIC, MMD, SMD)

\end{itemize}

\end{Details}
%
\begin{Value}
A transformed data frame with cleaned row names, reordered rows
(summary statistics at bottom), and the specified column names.
\end{Value}
%
\begin{SeeAlso}
\code{\LinkA{Tables2Tex}{Tables2Tex}} which calls this function to process full tables,
\code{\LinkA{GetTableEntry}{GetTableEntry}} for extracting table entries from models
\end{SeeAlso}
%
\begin{Examples}
\begin{ExampleCode}
## Not run: 
# Create a simple table
t_raw <- data.frame(Model1 = c("0.5 (2.1)*", "0.95", "100"))
row.names(t_raw) <- c("gdpPerCapita", "Adjusted R-squared", "Observations")

# Clean and rename
name_mat <- matrix(c("gdpPerCapita", "GDP per Capita"), ncol = 2)
t_clean <- FullTransformer(t_raw, COLNAMES_VEC = "Model 1",
                           name_conversion_matrix = name_mat)

## End(Not run)

\end{ExampleCode}
\end{Examples}
\HeaderA{GetTableEntry}{Extract regression results as a formatted table row}{GetTableEntry}
%
\begin{Description}
Builds a one-row data frame containing coefficient estimates, standard errors
or t-statistics, significance stars, and summary statistics from a fitted model.
Supports both analytical standard errors (including clustered and
heteroskedasticity-robust) and bootstrap-based inference.
\end{Description}
%
\begin{Usage}
\begin{verbatim}
GetTableEntry(
  my_lm,
  clust_id,
  iv_round = 2,
  NAME = "",
  iv = F,
  inParens = "tstat",
  seType = "analytical",
  bootstrap_source = "auto",
  bootstrap_reps = 999L,
  bootstrap_seed = NULL,
  bootstrap_dir = NULL,
  bootstrap_prefix = NULL,
  bootDataLocation = "./",
  bootDataNameTag = "Data",
  bootFactorVars = NULL,
  bootExcludeCovars = NULL,
  keepCoef1 = FALSE,
  superunit_covariateName = "country",
  superunit_label = "Countries"
)
\end{verbatim}
\end{Usage}
%
\begin{Arguments}
\begin{ldescription}
\item[\code{my\_lm}] A fitted model object (typically from \code{lm()}, \code{glm()},
or \code{ivreg()}). The data used to fit the model must be a data.frame.

\item[\code{clust\_id}] Character string giving the name of the clustering variable
for clustered standard errors. If \code{NULL}, heteroskedasticity-robust
standard errors (HC1) are used instead.

\item[\code{iv\_round}] Integer specifying the number of decimal places for rounding
estimates and statistics. Default is 2.

\item[\code{NAME}] Character string used as the column name for the resulting row.
Default is an empty string.

\item[\code{iv}] Logical; if \code{TRUE}, includes instrumental variable diagnostic
statistics (Weak instruments and Wu-Hausman tests) in the output. Only
applicable for models fitted with \code{ivreg()}. Default is \code{FALSE}.

\item[\code{inParens}] Character string specifying what to display in parentheses:
\code{"tstat"} (default) for t-statistics or \code{"se"} for standard errors.

\item[\code{seType}] Character string specifying the type of standard errors:
\code{"analytical"} (default) uses sandwich estimators, \code{"boot"} uses
bootstrap standard errors from either in-memory resampling or pre-computed
replication datasets.

\item[\code{bootstrap\_source}] Character string specifying how bootstrap replications
are sourced: \code{"auto"} (default), \code{"memory"}, or \code{"files"}.
Auto mode uses file-based bootstrap when \code{bootstrap\_dir} or
\code{bootstrap\_prefix} is supplied, and in-memory resampling otherwise.

\item[\code{bootstrap\_reps}] Integer giving the number of bootstrap replications to
use. For file-based bootstrap, this is the maximum number of replicate
files to load.

\item[\code{bootstrap\_seed}] Optional integer seed for reproducible in-memory
bootstrap resampling. Ignored for file-based bootstrap.

\item[\code{bootstrap\_dir}] Character string giving the folder path containing
bootstrap replication datasets when \code{bootstrap\_source = "files"}.

\item[\code{bootstrap\_prefix}] Character string giving the file name prefix for
bootstrap data files. Files should be named
\code{\{bootstrap\_prefix\}\_0.csv} for optional validation data and
\code{\{bootstrap\_prefix\}\_1.csv}, etc. for replications.

\item[\code{bootDataLocation}] Deprecated compatibility alias for
\code{bootstrap\_dir}.

\item[\code{bootDataNameTag}] Deprecated compatibility alias for
\code{bootstrap\_prefix}.

\item[\code{bootFactorVars}] Character vector of variable names to treat as factors
during bootstrap estimation. Deprecated and ignored in the new bootstrap
workflow.

\item[\code{bootExcludeCovars}] Character vector of covariate names to exclude from
the imputation step during bootstrap processing. Deprecated and ignored in
the new bootstrap workflow.

\item[\code{keepCoef1}] Logical; if \code{TRUE}, includes the first coefficient
(typically the intercept) in the output. Default is \code{FALSE}.

\item[\code{superunit\_covariateName}] Character string giving the name of the variable
used to count higher-level units (e.g., countries in panel data). Default
is \code{"country"}.

\item[\code{superunit\_label}] Character string used as the row label for the
higher-level unit count in the output table. Default is \code{"Countries"}.
\end{ldescription}
\end{Arguments}
%
\begin{Details}
The output format is designed for downstream processing by \code{\LinkA{Tables2Tex}{Tables2Tex}},
with coefficients formatted as "estimate (stat)*" where stat is either the
t-statistic or standard error (controlled by \code{inParens}), and * indicates
p < 0.05.
\end{Details}
%
\begin{Value}
A one-row data frame where each column corresponds to a coefficient
or summary statistic. Columns include:
\begin{itemize}

\item{} Formatted coefficient estimates with significance stars
\item{} Fit statistic (Adjusted R-squared for linear models, AIC for GLMs)
\item{} Number of observations
\item{} Count of higher-level units (e.g., countries)
\item{} IV diagnostics (if \code{iv = TRUE})

\end{itemize}

\end{Value}
%
\begin{SeeAlso}
\code{\LinkA{Tables2Tex}{Tables2Tex}} for generating complete LaTeX tables from
multiple models, \code{\LinkA{vcovCluster}{vcovCluster}} for the clustered standard error
implementation
\end{SeeAlso}
%
\begin{Examples}
\begin{ExampleCode}
## Not run: 
# Fit a simple linear model
data(mtcars)
fit <- lm(mpg ~ wt + hp, data = mtcars)

# Extract with robust standard errors
entry <- GetTableEntry(fit, clust_id = NULL, NAME = "Model 1")

# Extract with clustered standard errors
entry_clust <- GetTableEntry(fit, clust_id = "cyl", NAME = "Model 2")

## End(Not run)

\end{ExampleCode}
\end{Examples}
\HeaderA{heatMap}{Create an interpolated heat map from scattered data}{heatMap}
%
\begin{Description}
Interpolates irregularly-spaced \code{(x, y, z)} observations onto a regular
grid using the \pkg{interp} package and visualizes the result using
\code{\LinkA{image.plot}{image.plot}} from the \pkg{fields} package. This is useful
for visualizing relationships in data where observations are not on a regular
grid.
\end{Description}
%
\begin{Usage}
\begin{verbatim}
heatMap(
  x,
  y,
  z,
  main = "",
  N,
  yaxt = NULL,
  xlab = "",
  ylab = "",
  horizontal = F,
  useLog = "",
  legend.width = 1,
  ylim = NULL,
  xlim = NULL,
  zlim = NULL,
  add.legend = T,
  legend.only = F,
  vline = NULL,
  col_vline = "black",
  hline = NULL,
  col_hline = "black",
  cex.lab = 2,
  cex.main = 2,
  myCol = NULL,
  includeMarginals = F,
  marginalJitterSD_x = 0.01,
  marginalJitterSD_y = 0.01,
  openBrowser = F
)
\end{verbatim}
\end{Usage}
%
\begin{Arguments}
\begin{ldescription}
\item[\code{x}, \code{y}, \code{z}] Numeric vectors of equal length defining the coordinates and
values to interpolate. \code{x} and \code{y} are the spatial coordinates,
\code{z} is the value at each location.

\item[\code{main}] Character string for the plot title. Default is empty.

\item[\code{N}] Integer specifying the number of grid cells in each direction for
interpolation. Higher values produce smoother plots but take longer.

\item[\code{yaxt}] Optional character vector of labels for the y-axis. If \code{NULL}
(default), numeric axis labels are used.

\item[\code{xlab}, \code{ylab}] Character strings for axis labels. Default is empty.

\item[\code{horizontal}] Logical; if \code{TRUE}, draw the color legend horizontally
below the plot. Default is \code{FALSE} (vertical legend on right).

\item[\code{useLog}] Character string specifying which axes to log-transform.
Can include \code{"x"}, \code{"y"}, and/or \code{"z"} (e.g., \code{"xyz"}
for all axes). Default is empty (no log transformation).

\item[\code{legend.width}] Numeric value controlling the width of the color legend.
Default is 1.

\item[\code{xlim}, \code{ylim}, \code{zlim}] Numeric vectors of length 2 giving the plot limits for
each axis. If \code{NULL} (default), limits are computed from the data.

\item[\code{add.legend}] Logical; if \code{FALSE}, the color legend is suppressed.
Default is \code{TRUE}.

\item[\code{legend.only}] Logical; if \code{TRUE}, draw only the legend without the
main plot. Useful for creating custom layouts. Default is \code{FALSE}.

\item[\code{vline}, \code{hline}] Numeric values specifying positions of vertical or
horizontal reference lines. Default is \code{NULL} (no lines).

\item[\code{col\_vline}, \code{col\_hline}] Colors for the reference lines. Default is
\code{"black"}.

\item[\code{cex.lab}, \code{cex.main}] Numeric values for character expansion of axis labels
and title. Default is 2 for both.

\item[\code{myCol}] Optional color palette vector. If \code{NULL}, uses the default
\code{image.plot} palette.

\item[\code{includeMarginals}] Logical; if \code{TRUE}, adds 1D marginal rug plots
showing the distribution of x and y values. Default is \code{FALSE}.

\item[\code{marginalJitterSD\_x}, \code{marginalJitterSD\_y}] Numeric values controlling the
amount of jitter applied to marginal points, as a fraction of the standard
deviation. Default is 0.01.

\item[\code{openBrowser}] Logical; if \code{TRUE}, enters debug mode via
\code{browser()} for interactive inspection. Default is \code{FALSE}.
\end{ldescription}
\end{Arguments}
%
\begin{Value}
Invisibly returns \code{NULL}. Called for its side effect of producing
a plot.
\end{Value}
%
\begin{SeeAlso}
\code{\LinkA{heatmap2}{heatmap2}} for plotting matrices directly,
\code{\LinkA{MakeHeatMap}{MakeHeatMap}} for model-based prediction heat maps,
\code{\LinkA{image.plot}{image.plot}} for the underlying plotting function
\end{SeeAlso}
%
\begin{Examples}
\begin{ExampleCode}
## Not run: 
# Create some scattered data
set.seed(42)
x <- runif(100, 0, 10)
y <- runif(100, 0, 10)
z <- sin(x) * cos(y) + rnorm(100, sd = 0.1)

# Plot as interpolated heat map
heatMap(x, y, z, N = 50, main = "Example Heat Map",
        xlab = "X Variable", ylab = "Y Variable")

## End(Not run)

\end{ExampleCode}
\end{Examples}
\HeaderA{heatmap2}{Quick heat map visualization of a matrix}{heatmap2}
%
\begin{Description}
A convenience wrapper around \code{\LinkA{heatMap}{heatMap}} for fast visualization of
matrix-like objects. Can use either base R graphics or \pkg{ggplot2} for a
more polished appearance.
\end{Description}
%
\begin{Usage}
\begin{verbatim}
heatmap2(
  mat,
  row_labels = NULL,
  col_labels = NULL,
  use_gg = FALSE,
  scale = FALSE,
  log = FALSE,
  includeMarginals = FALSE,
  ...
)
\end{verbatim}
\end{Usage}
%
\begin{Arguments}
\begin{ldescription}
\item[\code{mat}] Numeric matrix or data frame to visualize. Will be coerced to a
matrix if necessary.

\item[\code{row\_labels}] Optional character vector of labels for rows. If \code{NULL}
(default), uses \code{rownames(mat)}.

\item[\code{col\_labels}] Optional character vector of labels for columns. If
\code{NULL} (default), uses \code{colnames(mat)}.

\item[\code{use\_gg}] Logical; if \code{TRUE}, creates a \pkg{ggplot2} plot instead
of using base graphics. Requires the \pkg{ggplot2} package. Default is
\code{FALSE}.

\item[\code{scale}] Logical; if \code{TRUE}, standardizes values (subtracts mean,
divides by SD) before plotting. Default is \code{FALSE}.

\item[\code{log}] Logical; if \code{TRUE}, applies natural log transformation to
values before plotting. Default is \code{FALSE}.

\item[\code{includeMarginals}] Logical; if \code{TRUE} and \code{use\_gg = TRUE},
adds marginal histograms using \pkg{ggExtra}. For base graphics, adds
marginal rug plots. Default is \code{FALSE}.

\item[\code{...}] Additional arguments passed to \code{\LinkA{heatMap}{heatMap}} when
\code{use\_gg = FALSE}.
\end{ldescription}
\end{Arguments}
%
\begin{Details}
When \code{use\_gg = TRUE}, the function creates a tile plot using
\code{\LinkA{geom\_tile}{geom.Rul.tile}} with a white-to-steelblue gradient. Marginal
histograms can be added using \pkg{ggExtra} if available.
\end{Details}
%
\begin{Value}
When \code{use\_gg = TRUE}, invisibly returns the ggplot object.
Otherwise, invisibly returns \code{NULL}. Called for its side effect of
producing a plot.
\end{Value}
%
\begin{SeeAlso}
\code{\LinkA{heatMap}{heatMap}} for the underlying scatter-to-grid
interpolation, \code{\LinkA{image2}{image2}} for simple matrix plotting
\end{SeeAlso}
%
\begin{Examples}
\begin{ExampleCode}
# Create a correlation matrix
cor_mat <- cor(mtcars[, 1:5])

# Quick base R plot
heatmap2(cor_mat, row_labels = colnames(cor_mat),
         col_labels = colnames(cor_mat))

## Not run: 
# ggplot2 version
heatmap2(cor_mat, use_gg = TRUE)

## End(Not run)

\end{ExampleCode}
\end{Examples}
\HeaderA{image2}{Plot a matrix as an image with intuitive orientation}{image2}
%
\begin{Description}
A wrapper around \code{\LinkA{image}{image}} that displays matrices with
the conventional orientation (first row at top, first column at left). The
base \code{image()} function displays matrices transposed and flipped, which
can be confusing. This function corrects that behavior.
\end{Description}
%
\begin{Usage}
\begin{verbatim}
image2(
  mat,
  xaxt = NULL,
  yaxt = NULL,
  col = NULL,
  main = NULL,
  scale_vec = c(1, 1.04),
  cex.axis = 1
)
\end{verbatim}
\end{Usage}
%
\begin{Arguments}
\begin{ldescription}
\item[\code{mat}] Numeric matrix to display as a heat map image.

\item[\code{xaxt}] Optional character vector of labels for the x-axis (columns).
If \code{NULL} (default), no x-axis labels are drawn.

\item[\code{yaxt}] Optional character vector of labels for the y-axis (rows).
If \code{NULL} (default), no y-axis labels are drawn.

\item[\code{col}] Optional vector of colors for the heat map. If \code{NULL},
uses the default \code{image()} color palette.

\item[\code{main}] Optional character string for the plot title.

\item[\code{scale\_vec}] Numeric vector of length 2 controlling the placement scaling
of tick labels on the x and y axes. Default is \code{c(1, 1.04)}.

\item[\code{cex.axis}] Numeric value for character expansion of axis tick labels.
Default is 1.
\end{ldescription}
\end{Arguments}
%
\begin{Value}
Invisibly returns \code{NULL}. Called for its side effect of producing
a plot.
\end{Value}
%
\begin{SeeAlso}
\code{\LinkA{heatmap2}{heatmap2}} for a higher-level heat map function,
\code{\LinkA{heatMap}{heatMap}} for interpolated heat maps from scattered data
\end{SeeAlso}
%
\begin{Examples}
\begin{ExampleCode}
# Create a simple matrix
mat <- matrix(1:12, nrow = 3, ncol = 4)
rownames(mat) <- c("A", "B", "C")
colnames(mat) <- c("W", "X", "Y", "Z")

# Plot with row/column labels
image2(mat, xaxt = colnames(mat), yaxt = rownames(mat), main = "Example")

\end{ExampleCode}
\end{Examples}
\HeaderA{MakeHeatMap}{Visualize two-way predictor effects as a heat map}{MakeHeatMap}
%
\begin{Description}
Creates a prediction surface showing how the expected outcome varies across
a grid of two predictor variables, holding all other predictors at their
mean (numeric) or mode (categorical) values. The result is saved as a PDF
file.
\end{Description}
%
\begin{Usage}
\begin{verbatim}
MakeHeatMap(
  factor1,
  factor2,
  outcome,
  dat,
  lm_obj,
  pdf_path,
  extrap_factor1 = 1,
  extrap_factor2 = 1,
  useLog = "",
  OUTCOME_SCALER = 1,
  OutcomeTransformFxn = function(x) {
     x
 },
  openBrowser = F
)
\end{verbatim}
\end{Usage}
%
\begin{Arguments}
\begin{ldescription}
\item[\code{factor1}, \code{factor2}] Character strings giving the names of the two
predictor variables to vary. These must be columns in \code{dat}.

\item[\code{outcome}] Character string giving the name of the outcome variable.
Used only for excluding from the summary computation.

\item[\code{dat}] Data frame containing the data used to fit the model. Used to
determine variable ranges and compute summary values.

\item[\code{lm\_obj}] Fitted linear model object (from \code{lm()} or similar) used
for generating predictions.

\item[\code{pdf\_path}] Character string giving the file path where the PDF plot
will be saved.

\item[\code{extrap\_factor1}, \code{extrap\_factor2}] Numeric multipliers controlling how far
beyond the observed data range the prediction grid extends. Values > 1
extrapolate beyond the data; values < 1 restrict to a subset of the range.
Default is 1 (exact data range).

\item[\code{useLog}] Character string specifying which axes to log-transform.
Can include \code{"x"}, \code{"y"}, and/or \code{"z"}. Default is empty
(no transformation).

\item[\code{OUTCOME\_SCALER}] Numeric value by which predictions are multiplied
before plotting. Useful for unit conversions. Default is 1.

\item[\code{OutcomeTransformFxn}] Function applied to predictions before scaling.
Default is the identity function \code{function(x) x}.

\item[\code{openBrowser}] Logical; if \code{TRUE}, pauses execution via
\code{browser()} for interactive debugging. Default is \code{FALSE}.
\end{ldescription}
\end{Arguments}
%
\begin{Details}
The function uses \code{\LinkA{colSummmary}{colSummmary}} to compute representative values
for predictors not being varied, then generates predictions across a 50x50
grid spanning the ranges of the two focal predictors. The predictions are
visualized using \code{\LinkA{heatMap}{heatMap}} with the \pkg{viridis} color palette.
\end{Details}
%
\begin{Value}
Invisibly returns \code{NULL}. The function is called for its side
effect of writing a PDF file to \code{pdf\_path}.
\end{Value}
%
\begin{SeeAlso}
\code{\LinkA{heatMap}{heatMap}} for the underlying plotting function,
\code{\LinkA{colSummmary}{colSummmary}} for computing representative predictor values
\end{SeeAlso}
%
\begin{Examples}
\begin{ExampleCode}
## Not run: 
# Fit a model
fit <- lm(mpg ~ wt + hp + cyl, data = mtcars)

# Create heat map of mpg as function of wt and hp
MakeHeatMap(
  factor1 = "wt",
  factor2 = "hp",
  outcome = "mpg",
  dat = mtcars,
  lm_obj = fit,
  pdf_path = "mpg_heatmap.pdf"
)

## End(Not run)

\end{ExampleCode}
\end{Examples}
\HeaderA{Stargazer2FullTable}{Convert a stargazer table to a self-contained longtable}{Stargazer2FullTable}
%
\begin{Description}
Transforms the LaTeX output from \code{\LinkA{stargazer}{stargazer}} into a
longtable environment suitable for multi-page regression tables. This function
handles the conversion automatically, including header repetition on continued
pages.
\end{Description}
%
\begin{Usage}
\begin{verbatim}
Stargazer2FullTable(stargazer_text, fontsize = "footnotesize")
\end{verbatim}
\end{Usage}
%
\begin{Arguments}
\begin{ldescription}
\item[\code{stargazer\_text}] Character vector containing the raw lines produced by
\code{stargazer()} (typically captured with \code{capture.output()}).

\item[\code{fontsize}] Character string specifying the LaTeX font size environment
to wrap the table in. Common values include \code{"footnotesize"},
\code{"scriptsize"}, or \code{"tiny"}. Default is \code{"footnotesize"}.
\end{ldescription}
\end{Arguments}
%
\begin{Details}
The function performs several transformations:
\begin{itemize}

\item{} Replaces \code{table} with \code{longtable} environment
\item{} Wraps content in the specified font size environment
\item{} Adds \code{\bsl{}endfirsthead} and \code{\bsl{}endhead} markers for header
repetition
\item{} Includes "(Continued from previous page)" note on subsequent pages
\item{} Removes incompatible stargazer formatting directives

\end{itemize}

\end{Details}
%
\begin{Value}
A character vector of LaTeX code ready to be written to a .tex file
using \code{write()} or \code{cat()}.
\end{Value}
%
\begin{Section}{LaTeX Requirements}

The output requires the \code{longtable} package in your LaTeX preamble:
\begin{alltt}\bsl{}usepackage\{longtable\}\end{alltt}

\end{Section}
%
\begin{SeeAlso}
\code{\LinkA{Tables2Tex}{Tables2Tex}} which calls this function internally,
\code{\LinkA{FullTransformer}{FullTransformer}} for table cleaning before conversion
\end{SeeAlso}
%
\begin{Examples}
\begin{ExampleCode}
## Not run: 
# Capture stargazer output and convert to longtable
library(stargazer)
fit <- lm(mpg ~ wt + hp, data = mtcars)
sg_output <- capture.output(stargazer(fit))

# Convert to longtable format
lt_output <- Stargazer2FullTable(sg_output, fontsize = "scriptsize")

# Write to file
write(lt_output, file = "full_table.tex")

## End(Not run)

\end{ExampleCode}
\end{Examples}
\HeaderA{Tables2Tex}{Generate publication-ready regression tables in LaTeX}{Tables2Tex}
%
\begin{Description}
Collates a list of regression model objects, extracts their coefficient
estimates and statistics via \code{\LinkA{GetTableEntry}{GetTableEntry}}, and writes
publication-ready LaTeX table files using \code{\LinkA{stargazer}{stargazer}}.
\end{Description}
%
\begin{Usage}
\begin{verbatim}
Tables2Tex(
  reg_list,
  clust_id,
  seType = "analytical",
  bootstrap_source = "auto",
  bootstrap_reps = 999L,
  bootstrap_seed = NULL,
  bootstrap_dir = NULL,
  bootstrap_prefix = NULL,
  checkmark_list = NULL,
  addrow_list = NULL,
  saveFolder = "./",
  nameTag = "Table",
  saveFull = T,
  tabCaption = "",
  model.names = NULL,
  NameConversionMat = NULL,
  DoFullTableKey = T,
  superunit_covariateName = "country",
  superunit_label = "Countries",
  font.size = "footnotesize",
  inParens = "tstat",
  keepCoef1 = FALSE,
  font.size.full = "footnotesize"
)
\end{verbatim}
\end{Usage}
%
\begin{Arguments}
\begin{ldescription}
\item[\code{reg\_list}] A list of fitted model objects, or a character vector of
object names to evaluate. Models should be fitted using \code{lm()},
\code{glm()}, or \code{ivreg()}. The data argument to these models
must be a data.frame (not a matrix or tibble).

\item[\code{clust\_id}] Character string giving the name of the clustering variable
for clustered standard errors, or \code{NULL} for heteroskedasticity-robust
standard errors.

\item[\code{seType}] Character string specifying the type of standard errors:
\code{"analytical"} (default) or \code{"boot"} for bootstrap.

\item[\code{bootstrap\_source}] Character string specifying how bootstrap replications
are sourced when \code{seType = "boot"}: \code{"auto"} (default),
\code{"memory"}, or \code{"files"}.

\item[\code{bootstrap\_reps}] Integer giving the number of bootstrap replications to
run in memory mode, or the maximum number of replicate files to load in
file mode.

\item[\code{bootstrap\_seed}] Optional integer seed for reproducible in-memory
bootstrap resampling.

\item[\code{bootstrap\_dir}] Character string giving the folder path containing
bootstrap replication datasets when \code{bootstrap\_source = "files"}.

\item[\code{bootstrap\_prefix}] Character string giving the file name prefix for
bootstrap data files, e.g. \code{\{bootstrap\_prefix\}\_1.csv}.

\item[\code{checkmark\_list}] Optional named list of binary vectors indicating which
models include certain features. Names become row labels, values of 1
produce checkmarks in the table.

\item[\code{addrow\_list}] Optional named list of vectors to append as additional
rows to the table. Each list element becomes a row with the element name
as the row label.

\item[\code{saveFolder}] Character string giving the folder path where LaTeX files
will be written. Default is the current directory.

\item[\code{nameTag}] Character string used as the base name for output files.
Files are named \code{tab\{nameTag\}\_SE\{seType\}.tex} and
\code{FULL\_tab\{nameTag\}\_SE\{seType\}.tex}.

\item[\code{saveFull}] Logical; if \code{TRUE} (default), also produces a full table
containing all coefficients using the longtable environment.

\item[\code{tabCaption}] Character string for the table caption.

\item[\code{model.names}] Optional character vector of column headings for each model.
If \code{NULL}, defaults to "Model 1", "Model 2", etc.

\item[\code{NameConversionMat}] Optional two-column matrix for filtering and renaming
row labels. Column 1 contains regex patterns to match, column 2 contains
replacement names. Rows not matching any pattern are dropped from the
condensed table (but retained in the full table).

\item[\code{DoFullTableKey}] Logical; if \code{TRUE} (default), adds a reference to
the full table in the condensed table's caption.

\item[\code{superunit\_covariateName}] Character string giving the variable name used
to count higher-level units (e.g., \code{"country"} for panel data).

\item[\code{superunit\_label}] Character string for the row label showing the count
of higher-level units. Default is \code{"Countries"}.

\item[\code{font.size}] Character string specifying the LaTeX font size for the
condensed table (e.g., \code{"footnotesize"}, \code{"scriptsize"}).

\item[\code{inParens}] Character string specifying what to display in parentheses:
\code{"tstat"} (default) or \code{"se"} for standard errors.

\item[\code{keepCoef1}] Logical; if \code{TRUE}, includes the first coefficient
(typically the intercept) in the output. Default is \code{FALSE}.

\item[\code{font.size.full}] Character string specifying the LaTeX font size for
the full table. Default is \code{"footnotesize"}.
\end{ldescription}
\end{Arguments}
%
\begin{Details}
This is the main user-facing function for the regression table workflow.
It produces two outputs:
\begin{enumerate}

\item{} A condensed table showing only the covariates specified in
\code{NameConversionMat} (or all if \code{NULL})
\item{} A full table (if \code{saveFull = TRUE}) containing all coefficients
in longtable format for multi-page output

\end{enumerate}

\end{Details}
%
\begin{Value}
Invisibly returns \code{NULL}. The function is called for its side
effect of writing LaTeX files to \code{saveFolder}.
\end{Value}
%
\begin{Section}{LaTeX Requirements}

The generated tables require the following LaTeX packages:
\begin{itemize}

\item{} \code{longtable} for full tables
\item{} \code{amssymb} for checkmark symbols

\end{itemize}

\end{Section}
%
\begin{SeeAlso}
\code{\LinkA{GetTableEntry}{GetTableEntry}} for extracting individual model results,
\code{\LinkA{FullTransformer}{FullTransformer}} for table cleaning,
\code{\LinkA{Stargazer2FullTable}{Stargazer2FullTable}} for longtable conversion
\end{SeeAlso}
%
\begin{Examples}
\begin{ExampleCode}
## Not run: 
# Fit multiple models
fit1 <- lm(mpg ~ wt, data = mtcars)
fit2 <- lm(mpg ~ wt + hp, data = mtcars)
fit3 <- lm(mpg ~ wt + hp + cyl, data = mtcars)

# Generate tables
Tables2Tex(
  reg_list = list(fit1, fit2, fit3),
  clust_id = NULL,
  saveFolder = "./tables/",
  nameTag = "MPG_Models",
  tabCaption = "Determinants of Fuel Efficiency",
  model.names = c("Base", "Controls", "Full")
)

## End(Not run)

\end{ExampleCode}
\end{Examples}
\HeaderA{vcovCluster}{Cluster-robust covariance matrix estimator}{vcovCluster}
%
\begin{Description}
Computes a clustered sandwich covariance matrix for a fitted model using the
method of Arellano (1987). This provides standard errors that are robust to
arbitrary within-cluster correlation.
\end{Description}
%
\begin{Usage}
\begin{verbatim}
vcovCluster(fm, clvar)
\end{verbatim}
\end{Usage}
%
\begin{Arguments}
\begin{ldescription}
\item[\code{fm}] A fitted model object (typically from \code{lm()} or \code{glm()}).
For \code{polr} models from \pkg{MASS}, predictions are handled specially.

\item[\code{clvar}] Character string giving the name of the clustering variable.
This variable must exist in the data used to fit \code{fm}.
\end{ldescription}
\end{Arguments}
%
\begin{Details}
The function applies the degrees-of-freedom correction
\eqn{(M/(M-1)) \times ((N-1)/(N-K))}{} where M is
the number of clusters, N is the number of observations, and K is the number
of parameters.
\end{Details}
%
\begin{Value}
A K x K covariance matrix where K is the number of model coefficients.
\end{Value}
%
\begin{References}
Arellano, M. (1987). Computing Robust Standard Errors for Within-Groups
Estimators. \emph{Oxford Bulletin of Economics and Statistics}, 49(4), 431-434.
\end{References}
%
\begin{SeeAlso}
\code{\LinkA{GetTableEntry}{GetTableEntry}} which uses this function for clustered
standard errors, \code{\LinkA{vcovHC}{vcovHC}} for heteroskedasticity-robust
covariance without clustering
\end{SeeAlso}
%
\begin{Examples}
\begin{ExampleCode}
## Not run: 
data(mtcars)
fit <- lm(mpg ~ wt + hp, data = mtcars)

# Cluster by number of cylinders
V_clust <- vcovCluster(fit, "cyl")

# Use with coeftest for clustered standard errors
library(lmtest)
coeftest(fit, vcov. = V_clust)

## End(Not run)

\end{ExampleCode}
\end{Examples}
\HeaderA{WidenMargins}{Widen margins in a LaTeX table string}{WidenMargins}
%
\begin{Description}
Adds horizontal margin adjustments to a LaTeX table by wrapping the table
content in an \code{adjustwidth} environment from the \pkg{ragged2e} package.
This is useful when tables are too wide to fit within standard margins.
\end{Description}
%
\begin{Usage}
\begin{verbatim}
WidenMargins(x)
\end{verbatim}
\end{Usage}
%
\begin{Arguments}
\begin{ldescription}
\item[\code{x}] Character vector containing the LaTeX table code, typically the
output from \code{\LinkA{stargazer}{stargazer}} or \code{\LinkA{Tables2Tex}{Tables2Tex}}.
\end{ldescription}
\end{Arguments}
%
\begin{Details}
The function extends margins by 0.5 inches on each side, effectively allowing
the table to be 1 inch wider than the text width.
\end{Details}
%
\begin{Value}
A character vector with the modified LaTeX code. The \code{adjustwidth}
environment is inserted inside the \code{table} environment.
\end{Value}
%
\begin{Section}{LaTeX Requirements}

The output requires the \pkg{ragged2e} package in your LaTeX preamble:
\begin{alltt}\bsl{}usepackage\{ragged2e\}\end{alltt}

\end{Section}
%
\begin{SeeAlso}
\code{\LinkA{Tables2Tex}{Tables2Tex}} for generating LaTeX tables
\end{SeeAlso}
%
\begin{Examples}
\begin{ExampleCode}
## Not run: 
# Read an existing table and widen its margins
tex_lines <- readLines("my_table.tex")
tex_wide <- WidenMargins(tex_lines)
writeLines(tex_wide, "my_table_wide.tex")

## End(Not run)

\end{ExampleCode}
\end{Examples}
\printindex{}
\end{document}
